\chapter{2006:Roger D. Kornberg}

\label{chap:chemistry2006}

The Nobel Prize in Chemistry for the year 2006 was awarded to Roger D. Kornberg.

\textbf{Motivation:}

for his studies of the molecular basis of eukaryotic transcription

\section{Roger D. Kornberg}

\subsection*{Early Life and Education}

Roger D. Kornberg was born in 1947 in St. Louis, Missouri, USA.

\subsection*{Career and Major Research}

Kornberg earned a bachelor's degree from Harvard University and a doctorate from Stanford University in 1972. After postdoctoral work at the MRC Laboratory of Molecular Biology in Cambridge with Francis Crick, he joined the faculty of Stanford University, where he remains. His father, Arthur Kornberg, received the Nobel Prize in Physiology or Medicine in 1959.

\subsection*{Nobel Prize-winning Work}

The work for which Roger D. Kornberg was awarded the Nobel Prize in Chemistry in 2006 was described by the Nobel Committee as: "for his studies of the molecular basis of eukaryotic transcription".

Kornberg determined the three-dimensional structure of RNA polymerase II, the enzyme that copies DNA into messenger RNA in the cells of plants, animals, and fungi, and showed how it works together with transcription factors to read the genetic information.

\subsection*{Significance and Impact}

His structural work revealed at atomic resolution how the transcription machinery selects and reads genes and how gene expression is regulated. It explained why transcription in complex organisms is controlled at the molecular level.

\subsection*{Later Career and Honors}

Kornberg continues to direct research on transcription and chromatin at the Stanford University School of Medicine.

\subsection*{Legacy}

The structure of RNA polymerase II stands as a foundation of modern molecular biology, clarifying how genetic information is expressed in all higher organisms.

\subsection*{Modern Developments}

Kornberg's X-ray crystallographic structure of RNA polymerase II revealed how DNA is transcribed into RNA, the first step in gene expression. His structure, showing the enzyme at work in an intact transcription complex, underpins the modern understanding of gene regulation and the design of transcription-targeting drugs.

\subsection*{Selected Publications}

"Structural Basis of Transcription: An RNA Polymerase II Elongation Complex at 3.3 Å Resolution" (Science, 2001)

\clearpage

\section*{Chapter Summary}

The 2006 prize honored work that revealed, at atomic resolution, how RNA polymerase II transcribes DNA into RNA in eukaryotic cells. It established the molecular basis of transcription and its regulation.
